\documentclass[11pt,a4paper]{article}
\usepackage[utf8]{inputenc}
\usepackage[T1]{fontenc}
\usepackage[spanish]{babel}
\usepackage{geometry,booktabs,longtable,array,xcolor,hyperref,amsmath}
\geometry{margin=1.7cm}
\hypersetup{colorlinks=true,urlcolor=blue}

\title{Metrados y Presupuesto Consolidado de Instalaciones Eléctricas\\Vivienda Unifamiliar de 2 Pisos - Proyecto Aquiles}
\author{Aquiles Taylor Ramos Yapo\\Universidad Nacional del Altiplano -- Ingeniería Mecánica Eléctrica}
\date{08 de junio de 2026}

\begin{document}
\maketitle

\section{INTRODUCCIÓN Y METODOLOGÍA}
El presente documento técnico detalla los metrados de materiales eléctricos y el presupuesto consolidado para el proyecto de instalaciones eléctricas interiores de la vivienda unifamiliar de dos pisos de propiedad de \textbf{Aquiles Taylor Ramos Yapo}, ubicada en el distrito de San Miguel, provincia de San Román y departamento de Puno.

Para la elaboración del metrado, se han contabilizado y validado las partidas directamente desde los planos eléctricos CAD vigentes (v7 para el primer piso y v6 para el segundo piso). Se han considerado todos los componentes necesarios para una instalación segura y reglamentaria bajo el Código Nacional de Electricidad (CNE - Utilización), tales como cables de alimentación de $16\text{ mm}^2$, conductores de derivados de $2.5\text{ mm}^2$ (alumbrado y tomacorrientes), canalizaciones de PVC SAP pesadas, cajas de fierro galvanizado, protecciones termomagnéticas y diferenciales de $30\text{ mA}$, y el sistema de puesta a tierra (SPAT).

\section{RESUMEN DE COTIZACIÓN MULTI-PROVEEDOR}
Para obtener un presupuesto real y eficiente de adquisición de materiales, se utilizó una herramienta de cotización multi-proveedor que compara precios en tiempo real en las principales tiendas locales (Promart, Sodimac y Maestro). El sistema incorpora una capa de conversión a unidades comerciales reales (por ejemplo, convirtiendo los metros de cable de diseño a rollos de 100 m, y los tubos a unidades de 3 m), calculando el subtotal en base a empaques cerrados y considerando los sobrantes.

Los resultados generales del análisis son:
\begin{itemize}
  \item \textbf{Costo total recomendado mixto (Compra Real):} S/ 3,968.02.
  \item \textbf{Costo neto equivalente (Valores de diseño):} S/ 3,892.40.
  \item \textbf{Número de partidas de metrado cotizadas:} 68 partidas consolidables.
  \item \textbf{Distribución recomendada del pedido para optimizar costos:}
    \begin{itemize}
      \item \textbf{Sodimac / Maestro:} S/ 2,169.42 (incluye cables en rollos de 100 m, tomacorrientes y protecciones generales).
      \item \textbf{Promart:} S/ 1,798.60 (incluye tuberías de PVC, interruptores, cajas de derivación y accesorios de montaje).
    \end{itemize}
\end{itemize}

\clearpage
% Fragmento listo para: % Fragmento listo para: % Fragmento listo para: \input{../presupuesto/presupuesto_final_aquiles.tex}
% Requiere: booktabs, longtable, array, xcolor, hyperref.
\section{Metrados de materiales electricos}
Se validaron 68 partidas usando los DXF vigentes, sus JSON documentales y los calculos previos del BOM.
\scriptsize
\begin{longtable}{p{1.7cm}p{7.5cm}r c p{2.2cm}}
\toprule Codigo & Descripcion & Cant. & Und. & Circuito \\ \midrule
\endfirsthead\toprule Codigo & Descripcion & Cant. & Und. & Circuito \\ \midrule\endhead
BOM-001 & Cable TW THW 16 mm2 (alimentador trifasico) & 40 & m & ALIM-TG \\
BOM-002 & Tubo PVC SAP 40 mm (alimentador) & 13 & m & ALIM-TG \\
BOM-003 & ITM 2P-63A & 1 & und & ALIM-TG \\
BOM-004 & Tablero general para 8 circuitos & 1 & und & GENERAL \\
BOM-005 & Interruptor diferencial 2P-40A-30mA & 1 & und & GENERAL \\
BOM-006 & Cable TW THW 2.5 mm2 - C1 & 33 & m & C1 \\
BOM-007 & Tubo PVC SAP 20 mm - C1 & 16 & m & C1 \\
BOM-008 & ITM 2P-10A - C1 & 1 & und & C1 \\
BOM-010 & Cable TW THW 2.5 mm2 - C2 & 40 & m & C2 \\
BOM-011 & Tubo PVC SAP 20 mm - C2 & 20 & m & C2 \\
BOM-012 & ITM 2P-10A - C2 & 1 & und & C2 \\
BOM-013 & Diferencial 2P-25A-30mA - C2 & 1 & und & C2 \\
BOM-015 & Cable TW THW 2.5 mm2 - C3 & 22 & m & C3 \\
BOM-016 & Tubo PVC SAP 20 mm - C3 & 11 & m & C3 \\
BOM-017 & ITM 2P-10A - C3 & 1 & und & C3 \\
BOM-018 & Diferencial 2P-25A-30mA - C3 & 1 & und & C3 \\
BOM-020 & Cable TW THW 2.5 mm2 - C4 & 48 & m & C4 \\
BOM-021 & Tubo PVC SAP 20 mm - C4 & 24 & m & C4 \\
BOM-022 & ITM 2P-10A - C4 & 1 & und & C4 \\
BOM-024 & Cable TW THW 2.5 mm2 - C5 & 55 & m & C5 \\
BOM-025 & Tubo PVC SAP 20 mm - C5 & 28 & m & C5 \\
BOM-026 & ITM 2P-10A - C5 & 1 & und & C5 \\
BOM-027 & Diferencial 2P-25A-30mA - C5 & 1 & und & C5 \\
BOM-029 & Cable TW THW 10 mm2 - C6 & 44 & m & C6 \\
BOM-030 & Tubo PVC SAP 20 mm - C6 & 22 & m & C6 \\
BOM-031 & ITM 2P-32A - C6 & 1 & und & C6 \\
BOM-032 & Diferencial 2P-40A-30mA - C6 & 1 & und & C6 \\
BOM-034 & Cable TW THW 2.5 mm2 - C7 & 33 & m & C7 \\
BOM-035 & Tubo PVC SAP 20 mm - C7 & 16 & m & C7 \\
BOM-036 & ITM 2P-10A - C7 & 1 & und & C7 \\
BOM-037 & Diferencial 2P-25A-30mA - C7 & 1 & und & C7 \\
BOM-039 & Cable TW THW 2.5 mm2 - C8 & 31 & m & C8 \\
BOM-040 & Tubo PVC SAP 20 mm - C8 & 15 & m & C8 \\
BOM-041 & ITM 2P-10A - C8 & 1 & und & C8 \\
BOM-042 & Diferencial 2P-25A-30mA - C8 & 1 & und & C8 \\
BOM-044 & Cable de tierra TW 10 mm2 (verde/amarillo) & 46 & m & SPAT \\
BOM-045 & Varilla de puesta a tierra 5/8 x 2.4 m & 1 & und & SPAT \\
BOM-046 & Cinta de aislar electrica & 3 & und & GENERAL \\
BOM-047 & Caja octogonal para luminaria - C1 & 4 & und & C1 \\
BOM-048 & Caja rectangular para interruptor - C1 & 4 & und & C1 \\
BOM-049 & Caja rectangular para tomacorriente - C2 & 7 & und & C2 \\
BOM-050 & Caja rectangular para tomacorriente especial - C3 & 1 & und & C3 \\
BOM-051 & Caja octogonal para luminaria - C4 & 11 & und & C4 \\
BOM-052 & Caja rectangular para interruptor - C4 & 7 & und & C4 \\
BOM-053 & Caja rectangular para tomacorriente - C5 & 12 & und & C5 \\
BOM-054 & Caja rectangular para salida especial - C6 & 3 & und & C6 \\
BOM-055 & Caja rectangular para tomacorriente especial - C7 & 1 & und & C7 \\
BOM-056 & Caja estanca IP55 para conexion de bomba - C8 & 1 & und & C8 \\
BOM-057 & Luminaria LED plafon 18 W - C1 & 4 & und & C1 \\
BOM-058 & Luminaria LED plafon 18 W - C4 & 11 & und & C4 \\
BOM-059 & Interruptor simple 10 A - C1 & 3 & und & C1 \\
BOM-060 & Interruptor conmutado 10 A - C1 & 1 & und & C1 \\
BOM-061 & Interruptor simple 10 A - C4 & 6 & und & C4 \\
BOM-062 & Interruptor conmutado 10 A - C4 & 1 & und & C4 \\
BOM-063 & Tomacorriente doble 2P+T 16 A - C2 & 7 & und & C2 \\
BOM-064 & Tomacorriente doble 2P+T 16 A - C5 & 12 & und & C5 \\
BOM-065 & Tomacorriente especial 2P+T 20 A - C3 & 1 & und & C3 \\
BOM-066 & Tomacorriente especial 2P+T 32 A - C6 & 3 & und & C6 \\
BOM-067 & Tomacorriente especial 2P+T 20 A - C7 & 1 & und & C7 \\
BOM-068 & Tablero T2 para 8 circuitos & 1 & und & T2 \\
BOM-069 & Caja de registro para puesta a tierra & 1 & und & SPAT \\
BOM-070 & Conector abrazadera para varilla de puesta a tierra & 1 & und & SPAT \\
BOM-071 & Union PVC SAP 20 mm & 46 & und & C1-C8 \\
BOM-072 & Curva PVC SAP 20 mm & 16 & und & C1-C8 \\
BOM-073 & Union PVC SAP 40 mm & 4 & und & ALIM-TG \\
BOM-074 & Cable TW THW 10 mm2 - alimentador TG a T2 & 75 & m & ALIM-T2 \\
BOM-075 & Tubo PVC SAP 32 mm - alimentador TG a T2 & 25 & m & ALIM-T2 \\
BOM-076 & ITM 2P-40A - alimentador T2 & 1 & und & ALIM-T2 \\
\bottomrule\end{longtable}\normalsize

\section{Cotizacion por proveedor}
\subsection{Grupo 1: Promart}
\scriptsize
\begin{longtable}{p{6.0cm}r r r r c}
\toprule Material & Metrado & Compra & P. comercial & Subtotal & Tipo \\ \midrule
\endfirsthead\toprule Material & Metrado & Compra & P. comercial & Subtotal & Tipo \\ \midrule\endhead
cinta aislante electrica & 3 und & 3 unidad & 2.00 & 6.00 & V \\
curva pvc sap electrica 20 mm & 16 und & 16 unidad & 2.20 & 35.20 & V \\
union pvc sap electrica 20 mm & 46 und & 46 unidad & 1.10 & 50.60 & V \\
union pvc sap electrica 40 mm & 4 und & 4 unidad & 2.50 & 10.00 & E \\
cable electrico thw 10 mm2 cobre & 119 m & 119 metro & 4.00 & 476.00 & E \\
cable electrico thw 10 mm2 cobre & 46 m & 46 metro & 3.50 & 161.00 & E \\
cable electrico thw 16 mm2 cobre & 40 m & 40 metro & 6.00 & 240.00 & E \\
cable electrico thw 2.5 mm2 cobre & 262 m & 3 rollo & 209.90 & 629.70 & E \\
caja electrica estanca ip55 & 1 und & 1 unidad & 18.00 & 18.00 & E \\
caja electrica octogonal & 15 und & 15 unidad & 12.00 & 180.00 & E \\
caja electrica rectangular & 35 und & 35 paquete & 1.94 & 67.90 & E \\
caja registro puesta a tierra & 1 und & 1 unidad & 35.00 & 35.00 & E \\
tubo pvc sap electrico 20 mm & 152 m & 51 tubo & 5.90 & 300.90 & V \\
tubo pvc sap electrico 32 mm & 25 m & 9 tubo & 19.90 & 179.10 & V \\
tubo pvc sap electrico 40 mm & 13 m & 5 tubo & 12.00 & 60.00 & E \\
interruptor diferencial 2 polos 25A 30mA & 5 und & 5 unidad & 49.49 & 247.47 & E \\
interruptor diferencial 2 polos 40A 30mA & 2 und & 2 unidad & 65.90 & 131.80 & E \\
interruptor conmutado 10 a & 2 und & 2 unidad & 13.90 & 27.80 & E \\
interruptor simple 10 a & 9 und & 9 unidad & 5.90 & 53.10 & E \\
interruptor termomagnetico 2 polos 10A & 7 und & 7 unidad & 38.99 & 272.93 & E \\
interruptor termomagnetico 2 polos 32A & 1 und & 1 unidad & 37.90 & 37.90 & V \\
interruptor termomagnetico 2 polos 40A & 1 und & 1 unidad & 34.99 & 34.99 & E \\
interruptor termomagnetico 2 polos 63A & 1 und & 1 unidad & 35.45 & 35.45 & E \\
luminaria led plafon 18 w & 15 und & 15 unidad & 20.00 & 300.00 & E \\
conector abrazadera para varilla de puesta a tierra & 1 und & 1 unidad & 150.00 & 150.00 & E \\
varilla puesta a tierra cobre 5/8 2.4 m & 1 und & 1 unidad & 25.00 & 25.00 & E \\
tablero electrico 8 polos empotrable & 2 und & 2 unidad & 40.90 & 81.80 & V \\
tomacorriente doble 2p+t 16 a & 19 und & 19 unidad & 37.21 & 706.99 & E \\
tomacorriente especial 2p+t 20 a & 2 und & 2 unidad & 22.00 & 44.00 & E \\
tomacorriente especial 2p+t 32 a & 3 und & 3 unidad & 35.00 & 105.00 & E \\
\bottomrule\end{longtable}\normalsize

\subsection{Grupo 2: Sodimac/Maestro}
\scriptsize
\begin{longtable}{p{6.0cm}r r r r c}
\toprule Material & Metrado & Compra & P. comercial & Subtotal & Tipo \\ \midrule
\endfirsthead\toprule Material & Metrado & Compra & P. comercial & Subtotal & Tipo \\ \midrule\endhead
cinta aislante electrica & 3 und & 3 unidad & 6.72 & 20.16 & V \\
curva pvc sap electrica 20 mm & 16 und & 16 unidad & 2.60 & 41.60 & V \\
union pvc sap electrica 20 mm & 46 und & 46 unidad & 1.10 & 50.60 & E \\
union pvc sap electrica 40 mm & 4 und & 4 unidad & 2.50 & 10.00 & E \\
cable electrico thw 10 mm2 cobre & 119 m & 119 metro & 4.00 & 476.00 & E \\
cable electrico thw 10 mm2 cobre & 46 m & 46 metro & 3.50 & 161.00 & E \\
cable electrico thw 16 mm2 cobre & 40 m & 40 m & 19.90 & 796.00 & E \\
cable electrico thw 2.5 mm2 cobre & 262 m & 3 rollo & 209.90 & 629.70 & V \\
caja electrica estanca ip55 & 1 und & 1 unidad & 18.00 & 18.00 & E \\
caja electrica octogonal & 15 und & 15 paquete & 1.90 & 28.50 & V \\
caja electrica rectangular & 35 und & 35 paquete & 1.90 & 66.50 & V \\
caja registro puesta a tierra & 1 und & 1 unidad & 35.00 & 35.00 & E \\
tubo pvc sap electrico 20 mm & 152 m & 51 tubo & 6.90 & 351.90 & V \\
tubo pvc sap electrico 32 mm & 25 m & 9 tubo & 17.51 & 157.59 & V \\
tubo pvc sap electrico 40 mm & 13 m & 5 tubo & 12.00 & 60.00 & E \\
interruptor diferencial 2 polos 25A 30mA & 5 und & 5 unidad & 49.00 & 245.00 & V \\
interruptor diferencial 2 polos 40A 30mA & 2 und & 2 unidad & 65.90 & 131.80 & V \\
interruptor conmutado 10 a & 2 und & 2 unidad & 21.90 & 43.80 & E \\
interruptor simple 10 a & 9 und & 9 unidad & 43.90 & 395.10 & E \\
interruptor termomagnetico 2 polos 10A & 7 und & 7 unidad & 38.99 & 272.93 & V \\
interruptor termomagnetico 2 polos 32A & 1 und & 1 unidad & 31.90 & 31.90 & V \\
interruptor termomagnetico 2 polos 40A & 1 und & 1 unidad & 34.99 & 34.99 & V \\
interruptor termomagnetico 2 polos 63A & 1 und & 1 unidad & 29.90 & 29.90 & V \\
luminaria led plafon 18 w & 15 und & 15 unidad & 29.90 & 448.50 & E \\
conector abrazadera para varilla de puesta a tierra & 1 und & 1 unidad & 18.00 & 18.00 & V \\
varilla puesta a tierra cobre 5/8 2.4 m & 1 und & 1 unidad & 322.00 & 322.00 & E \\
tablero electrico 8 polos empotrable & 2 und & 2 unidad & 33.90 & 67.80 & V \\
tomacorriente doble 2p+t 16 a & 19 und & 19 unidad & 19.49 & 370.31 & V \\
tomacorriente especial 2p+t 20 a & 2 und & 2 unidad & 16.90 & 33.80 & E \\
tomacorriente especial 2p+t 32 a & 3 und & 3 unidad & 16.90 & 50.70 & E \\
\bottomrule\end{longtable}\normalsize

\subsection{Grupo 3: Mercado Libre Peru}
\scriptsize
\begin{longtable}{p{6.0cm}r r r r c}
\toprule Material & Metrado & Compra & P. comercial & Subtotal & Tipo \\ \midrule
\endfirsthead\toprule Material & Metrado & Compra & P. comercial & Subtotal & Tipo \\ \midrule\endhead
cinta aislante electrica & 3 und & 3 unidad & 5.71 & 17.13 & E \\
curva pvc sap electrica 20 mm & 16 und & 16 unidad & 2.48 & 39.60 & E \\
union pvc sap electrica 20 mm & 46 und & 46 unidad & 1.10 & 50.60 & E \\
union pvc sap electrica 40 mm & 4 und & 4 unidad & 2.50 & 10.00 & E \\
cable electrico thw 10 mm2 cobre & 119 m & 119 metro & 4.00 & 476.00 & E \\
cable electrico thw 10 mm2 cobre & 46 m & 46 metro & 3.50 & 161.00 & E \\
cable electrico thw 16 mm2 cobre & 40 m & 40 metro & 6.00 & 240.00 & E \\
cable electrico thw 2.5 mm2 cobre & 262 m & 3 rollo & 209.90 & 629.70 & E \\
caja electrica estanca ip55 & 1 und & 1 unidad & 18.00 & 18.00 & E \\
caja electrica octogonal & 15 und & 15 paquete & 1.90 & 28.50 & E \\
caja electrica rectangular & 35 und & 35 paquete & 1.94 & 67.90 & E \\
caja registro puesta a tierra & 1 und & 1 unidad & 35.00 & 35.00 & E \\
tubo pvc sap electrico 20 mm & 152 m & 51 tubo & 8.40 & 428.40 & E \\
tubo pvc sap electrico 32 mm & 25 m & 9 tubo & 18.31 & 164.76 & E \\
tubo pvc sap electrico 40 mm & 13 m & 5 tubo & 12.00 & 60.00 & E \\
interruptor diferencial 2 polos 25A 30mA & 5 und & 5 unidad & 49.49 & 247.47 & E \\
interruptor diferencial 2 polos 40A 30mA & 2 und & 2 unidad & 65.90 & 131.80 & E \\
interruptor conmutado 10 a & 2 und & 2 unidad & 18.00 & 36.00 & E \\
interruptor simple 10 a & 9 und & 9 unidad & 12.00 & 108.00 & E \\
interruptor termomagnetico 2 polos 10A & 7 und & 7 unidad & 38.99 & 272.93 & E \\
interruptor termomagnetico 2 polos 32A & 1 und & 1 unidad & 35.17 & 35.17 & E \\
interruptor termomagnetico 2 polos 40A & 1 und & 1 unidad & 34.99 & 34.99 & E \\
interruptor termomagnetico 2 polos 63A & 1 und & 1 unidad & 35.45 & 35.45 & E \\
luminaria led plafon 18 w & 15 und & 15 unidad & 25.00 & 375.00 & E \\
conector abrazadera para varilla de puesta a tierra & 1 und & 1 unidad & 20.45 & 20.45 & E \\
varilla puesta a tierra cobre 5/8 2.4 m & 1 und & 1 unidad & 25.00 & 25.00 & E \\
tablero electrico 8 polos empotrable & 2 und & 2 unidad & 50.01 & 100.02 & E \\
tomacorriente doble 2p+t 16 a & 19 und & 19 unidad & 37.21 & 706.99 & E \\
tomacorriente especial 2p+t 20 a & 2 und & 2 unidad & 22.00 & 44.00 & E \\
tomacorriente especial 2p+t 32 a & 3 und & 3 unidad & 35.00 & 105.00 & E \\
\bottomrule\end{longtable}\normalsize

\section{Comparativa economica}
\scriptsize
\begin{longtable}{p{4.5cm}rrrrp{2.5cm}r}
\toprule Material & Promart & Sod./Mae. & M. Libre & Mejor & Proveedor & Recomendado \\ \midrule
\endfirsthead\toprule Material & Promart & Sod./Mae. & M. Libre & Mejor & Proveedor & Recomendado \\ \midrule\endhead
cable electrico thw 10 mm2 cobre & 476.00 & 476.00 & 476.00 & Promart & 476.00 \\
cable electrico thw 10 mm2 cobre & 161.00 & 161.00 & 161.00 & Promart & 161.00 \\
cable electrico thw 16 mm2 cobre & 240.00 & 796.00 & 240.00 & Promart & 240.00 \\
cable electrico thw 2.5 mm2 cobre & 629.70 & 629.70 & 629.70 & Sodimac/Maestro & 629.70 \\
caja electrica estanca ip55 & 18.00 & 18.00 & 18.00 & Promart & 18.00 \\
caja electrica octogonal & 180.00 & 28.50 & 28.50 & Sodimac/Maestro & 28.50 \\
caja electrica rectangular & 67.90 & 66.50 & 67.90 & Sodimac/Maestro & 66.50 \\
caja registro puesta a tierra & 35.00 & 35.00 & 35.00 & Promart & 35.00 \\
cinta aislante electrica & 6.00 & 20.16 & 17.13 & Promart & 6.00 \\
conector abrazadera para varilla de puesta a tierra & 150.00 & 18.00 & 20.45 & Sodimac/Maestro & 18.00 \\
curva pvc sap electrica 20 mm & 35.20 & 41.60 & 39.60 & Promart & 35.20 \\
interruptor conmutado 10 a & 27.80 & 43.80 & 36.00 & Promart & 27.80 \\
interruptor diferencial 2 polos 25A 30mA & 247.47 & 245.00 & 247.47 & Sodimac/Maestro & 245.00 \\
interruptor diferencial 2 polos 40A 30mA & 131.80 & 131.80 & 131.80 & Sodimac/Maestro & 131.80 \\
interruptor simple 10 a & 53.10 & 395.10 & 108.00 & Promart & 53.10 \\
interruptor termomagnetico 2 polos 10A & 272.93 & 272.93 & 272.93 & Sodimac/Maestro & 272.93 \\
interruptor termomagnetico 2 polos 32A & 37.90 & 31.90 & 35.17 & Sodimac/Maestro & 31.90 \\
interruptor termomagnetico 2 polos 40A & 34.99 & 34.99 & 34.99 & Sodimac/Maestro & 34.99 \\
interruptor termomagnetico 2 polos 63A & 35.45 & 29.90 & 35.45 & Sodimac/Maestro & 29.90 \\
luminaria led plafon 18 w & 300.00 & 448.50 & 375.00 & Promart & 300.00 \\
tablero electrico 8 polos empotrable & 81.80 & 67.80 & 100.02 & Sodimac/Maestro & 67.80 \\
tomacorriente doble 2p+t 16 a & 706.99 & 370.31 & 706.99 & Sodimac/Maestro & 370.31 \\
tomacorriente especial 2p+t 20 a & 44.00 & 33.80 & 44.00 & Sodimac/Maestro & 33.80 \\
tomacorriente especial 2p+t 32 a & 105.00 & 50.70 & 105.00 & Sodimac/Maestro & 50.70 \\
tubo pvc sap electrico 20 mm & 300.90 & 351.90 & 428.40 & Promart & 300.90 \\
tubo pvc sap electrico 32 mm & 179.10 & 157.59 & 164.76 & Sodimac/Maestro & 157.59 \\
tubo pvc sap electrico 40 mm & 60.00 & 60.00 & 60.00 & Promart & 60.00 \\
union pvc sap electrica 20 mm & 50.60 & 50.60 & 50.60 & Promart & 50.60 \\
union pvc sap electrica 40 mm & 10.00 & 10.00 & 10.00 & Promart & 10.00 \\
varilla puesta a tierra cobre 5/8 2.4 m & 25.00 & 322.00 & 25.00 & Promart & 25.00 \\
\bottomrule\end{longtable}\normalsize

\section{Presupuesto recomendado}
\scriptsize
\begin{longtable}{p{5.5cm}p{2.5cm}r r p{3.0cm}}
\toprule Material & Proveedor & Compra & Subtotal & Tipo de precio \\ \midrule
\endfirsthead\toprule Material & Proveedor & Compra & Subtotal & Tipo de precio \\ \midrule\endhead
cinta aislante electrica & Promart & 3 unidad & 6.00 & Verificado (Tienda) \\
curva pvc sap electrica 20 mm & Promart & 16 unidad & 35.20 & Verificado (Tienda) \\
union pvc sap electrica 20 mm & Promart & 46 unidad & 50.60 & Verificado (Tienda) \\
union pvc sap electrica 40 mm & Promart & 4 unidad & 10.00 & Estimado (Referencia) \\
cable electrico thw 10 mm2 cobre & Promart & 119 metro & 476.00 & Estimado (Referencia) \\
cable electrico thw 10 mm2 cobre & Promart & 46 metro & 161.00 & Estimado (Referencia) \\
cable electrico thw 16 mm2 cobre & Promart & 40 metro & 240.00 & Estimado (Referencia) \\
cable electrico thw 2.5 mm2 cobre & Sodimac/Maestro & 3 rollo & 629.70 & Verificado (Tienda) \\
caja electrica estanca ip55 & Promart & 1 unidad & 18.00 & Estimado (Referencia) \\
caja electrica octogonal & Sodimac/Maestro & 15 paquete & 28.50 & Verificado (Tienda) \\
caja electrica rectangular & Sodimac/Maestro & 35 paquete & 66.50 & Verificado (Tienda) \\
caja registro puesta a tierra & Promart & 1 unidad & 35.00 & Estimado (Referencia) \\
tubo pvc sap electrico 20 mm & Promart & 51 tubo & 300.90 & Verificado (Tienda) \\
tubo pvc sap electrico 32 mm & Sodimac/Maestro & 9 tubo & 157.59 & Verificado (Tienda) \\
tubo pvc sap electrico 40 mm & Promart & 5 tubo & 60.00 & Estimado (Referencia) \\
interruptor diferencial 2 polos 25A 30mA & Sodimac/Maestro & 5 unidad & 245.00 & Verificado (Tienda) \\
interruptor diferencial 2 polos 40A 30mA & Sodimac/Maestro & 2 unidad & 131.80 & Verificado (Tienda) \\
interruptor conmutado 10 a & Promart & 2 unidad & 27.80 & Estimado (Equivalente) \\
interruptor simple 10 a & Promart & 9 unidad & 53.10 & Estimado (Equivalente) \\
interruptor termomagnetico 2 polos 10A & Sodimac/Maestro & 7 unidad & 272.93 & Verificado (Tienda) \\
interruptor termomagnetico 2 polos 32A & Sodimac/Maestro & 1 unidad & 31.90 & Verificado (Tienda) \\
interruptor termomagnetico 2 polos 40A & Sodimac/Maestro & 1 unidad & 34.99 & Verificado (Tienda) \\
interruptor termomagnetico 2 polos 63A & Sodimac/Maestro & 1 unidad & 29.90 & Verificado (Tienda) \\
luminaria led plafon 18 w & Promart & 15 unidad & 300.00 & Estimado (Equivalente) \\
conector abrazadera para varilla de puesta a tierra & Sodimac/Maestro & 1 unidad & 18.00 & Verificado (Tienda) \\
varilla puesta a tierra cobre 5/8 2.4 m & Promart & 1 unidad & 25.00 & Estimado (Referencia) \\
tablero electrico 8 polos empotrable & Sodimac/Maestro & 2 unidad & 67.80 & Verificado (Tienda) \\
tomacorriente doble 2p+t 16 a & Sodimac/Maestro & 19 unidad & 370.31 & Verificado (Tienda) \\
tomacorriente especial 2p+t 20 a & Sodimac/Maestro & 2 unidad & 33.80 & Estimado (Equivalente) \\
tomacorriente especial 2p+t 32 a & Sodimac/Maestro & 3 unidad & 50.70 & Estimado (Equivalente) \\

\midrule\multicolumn{3}{r}{\textbf{TOTAL}} & \textbf{S/ 3,968.02} & \\ \bottomrule\end{longtable}\normalsize

% Requiere: booktabs, longtable, array, xcolor, hyperref.
\section{Metrados de materiales electricos}
Se validaron 68 partidas usando los DXF vigentes, sus JSON documentales y los calculos previos del BOM.
\scriptsize
\begin{longtable}{p{1.7cm}p{7.5cm}r c p{2.2cm}}
\toprule Codigo & Descripcion & Cant. & Und. & Circuito \\ \midrule
\endfirsthead\toprule Codigo & Descripcion & Cant. & Und. & Circuito \\ \midrule\endhead
BOM-001 & Cable TW THW 16 mm2 (alimentador trifasico) & 40 & m & ALIM-TG \\
BOM-002 & Tubo PVC SAP 40 mm (alimentador) & 13 & m & ALIM-TG \\
BOM-003 & ITM 2P-63A & 1 & und & ALIM-TG \\
BOM-004 & Tablero general para 8 circuitos & 1 & und & GENERAL \\
BOM-005 & Interruptor diferencial 2P-40A-30mA & 1 & und & GENERAL \\
BOM-006 & Cable TW THW 2.5 mm2 - C1 & 33 & m & C1 \\
BOM-007 & Tubo PVC SAP 20 mm - C1 & 16 & m & C1 \\
BOM-008 & ITM 2P-10A - C1 & 1 & und & C1 \\
BOM-010 & Cable TW THW 2.5 mm2 - C2 & 40 & m & C2 \\
BOM-011 & Tubo PVC SAP 20 mm - C2 & 20 & m & C2 \\
BOM-012 & ITM 2P-10A - C2 & 1 & und & C2 \\
BOM-013 & Diferencial 2P-25A-30mA - C2 & 1 & und & C2 \\
BOM-015 & Cable TW THW 2.5 mm2 - C3 & 22 & m & C3 \\
BOM-016 & Tubo PVC SAP 20 mm - C3 & 11 & m & C3 \\
BOM-017 & ITM 2P-10A - C3 & 1 & und & C3 \\
BOM-018 & Diferencial 2P-25A-30mA - C3 & 1 & und & C3 \\
BOM-020 & Cable TW THW 2.5 mm2 - C4 & 48 & m & C4 \\
BOM-021 & Tubo PVC SAP 20 mm - C4 & 24 & m & C4 \\
BOM-022 & ITM 2P-10A - C4 & 1 & und & C4 \\
BOM-024 & Cable TW THW 2.5 mm2 - C5 & 55 & m & C5 \\
BOM-025 & Tubo PVC SAP 20 mm - C5 & 28 & m & C5 \\
BOM-026 & ITM 2P-10A - C5 & 1 & und & C5 \\
BOM-027 & Diferencial 2P-25A-30mA - C5 & 1 & und & C5 \\
BOM-029 & Cable TW THW 10 mm2 - C6 & 44 & m & C6 \\
BOM-030 & Tubo PVC SAP 20 mm - C6 & 22 & m & C6 \\
BOM-031 & ITM 2P-32A - C6 & 1 & und & C6 \\
BOM-032 & Diferencial 2P-40A-30mA - C6 & 1 & und & C6 \\
BOM-034 & Cable TW THW 2.5 mm2 - C7 & 33 & m & C7 \\
BOM-035 & Tubo PVC SAP 20 mm - C7 & 16 & m & C7 \\
BOM-036 & ITM 2P-10A - C7 & 1 & und & C7 \\
BOM-037 & Diferencial 2P-25A-30mA - C7 & 1 & und & C7 \\
BOM-039 & Cable TW THW 2.5 mm2 - C8 & 31 & m & C8 \\
BOM-040 & Tubo PVC SAP 20 mm - C8 & 15 & m & C8 \\
BOM-041 & ITM 2P-10A - C8 & 1 & und & C8 \\
BOM-042 & Diferencial 2P-25A-30mA - C8 & 1 & und & C8 \\
BOM-044 & Cable de tierra TW 10 mm2 (verde/amarillo) & 46 & m & SPAT \\
BOM-045 & Varilla de puesta a tierra 5/8 x 2.4 m & 1 & und & SPAT \\
BOM-046 & Cinta de aislar electrica & 3 & und & GENERAL \\
BOM-047 & Caja octogonal para luminaria - C1 & 4 & und & C1 \\
BOM-048 & Caja rectangular para interruptor - C1 & 4 & und & C1 \\
BOM-049 & Caja rectangular para tomacorriente - C2 & 7 & und & C2 \\
BOM-050 & Caja rectangular para tomacorriente especial - C3 & 1 & und & C3 \\
BOM-051 & Caja octogonal para luminaria - C4 & 11 & und & C4 \\
BOM-052 & Caja rectangular para interruptor - C4 & 7 & und & C4 \\
BOM-053 & Caja rectangular para tomacorriente - C5 & 12 & und & C5 \\
BOM-054 & Caja rectangular para salida especial - C6 & 3 & und & C6 \\
BOM-055 & Caja rectangular para tomacorriente especial - C7 & 1 & und & C7 \\
BOM-056 & Caja estanca IP55 para conexion de bomba - C8 & 1 & und & C8 \\
BOM-057 & Luminaria LED plafon 18 W - C1 & 4 & und & C1 \\
BOM-058 & Luminaria LED plafon 18 W - C4 & 11 & und & C4 \\
BOM-059 & Interruptor simple 10 A - C1 & 3 & und & C1 \\
BOM-060 & Interruptor conmutado 10 A - C1 & 1 & und & C1 \\
BOM-061 & Interruptor simple 10 A - C4 & 6 & und & C4 \\
BOM-062 & Interruptor conmutado 10 A - C4 & 1 & und & C4 \\
BOM-063 & Tomacorriente doble 2P+T 16 A - C2 & 7 & und & C2 \\
BOM-064 & Tomacorriente doble 2P+T 16 A - C5 & 12 & und & C5 \\
BOM-065 & Tomacorriente especial 2P+T 20 A - C3 & 1 & und & C3 \\
BOM-066 & Tomacorriente especial 2P+T 32 A - C6 & 3 & und & C6 \\
BOM-067 & Tomacorriente especial 2P+T 20 A - C7 & 1 & und & C7 \\
BOM-068 & Tablero T2 para 8 circuitos & 1 & und & T2 \\
BOM-069 & Caja de registro para puesta a tierra & 1 & und & SPAT \\
BOM-070 & Conector abrazadera para varilla de puesta a tierra & 1 & und & SPAT \\
BOM-071 & Union PVC SAP 20 mm & 46 & und & C1-C8 \\
BOM-072 & Curva PVC SAP 20 mm & 16 & und & C1-C8 \\
BOM-073 & Union PVC SAP 40 mm & 4 & und & ALIM-TG \\
BOM-074 & Cable TW THW 10 mm2 - alimentador TG a T2 & 75 & m & ALIM-T2 \\
BOM-075 & Tubo PVC SAP 32 mm - alimentador TG a T2 & 25 & m & ALIM-T2 \\
BOM-076 & ITM 2P-40A - alimentador T2 & 1 & und & ALIM-T2 \\
\bottomrule\end{longtable}\normalsize

\section{Cotizacion por proveedor}
\subsection{Grupo 1: Promart}
\scriptsize
\begin{longtable}{p{6.0cm}r r r r c}
\toprule Material & Metrado & Compra & P. comercial & Subtotal & Tipo \\ \midrule
\endfirsthead\toprule Material & Metrado & Compra & P. comercial & Subtotal & Tipo \\ \midrule\endhead
cinta aislante electrica & 3 und & 3 unidad & 2.00 & 6.00 & V \\
curva pvc sap electrica 20 mm & 16 und & 16 unidad & 2.20 & 35.20 & V \\
union pvc sap electrica 20 mm & 46 und & 46 unidad & 1.10 & 50.60 & V \\
union pvc sap electrica 40 mm & 4 und & 4 unidad & 2.50 & 10.00 & E \\
cable electrico thw 10 mm2 cobre & 119 m & 119 metro & 4.00 & 476.00 & E \\
cable electrico thw 10 mm2 cobre & 46 m & 46 metro & 3.50 & 161.00 & E \\
cable electrico thw 16 mm2 cobre & 40 m & 40 metro & 6.00 & 240.00 & E \\
cable electrico thw 2.5 mm2 cobre & 262 m & 3 rollo & 209.90 & 629.70 & E \\
caja electrica estanca ip55 & 1 und & 1 unidad & 18.00 & 18.00 & E \\
caja electrica octogonal & 15 und & 15 unidad & 12.00 & 180.00 & E \\
caja electrica rectangular & 35 und & 35 paquete & 1.94 & 67.90 & E \\
caja registro puesta a tierra & 1 und & 1 unidad & 35.00 & 35.00 & E \\
tubo pvc sap electrico 20 mm & 152 m & 51 tubo & 5.90 & 300.90 & V \\
tubo pvc sap electrico 32 mm & 25 m & 9 tubo & 19.90 & 179.10 & V \\
tubo pvc sap electrico 40 mm & 13 m & 5 tubo & 12.00 & 60.00 & E \\
interruptor diferencial 2 polos 25A 30mA & 5 und & 5 unidad & 49.49 & 247.47 & E \\
interruptor diferencial 2 polos 40A 30mA & 2 und & 2 unidad & 65.90 & 131.80 & E \\
interruptor conmutado 10 a & 2 und & 2 unidad & 13.90 & 27.80 & E \\
interruptor simple 10 a & 9 und & 9 unidad & 5.90 & 53.10 & E \\
interruptor termomagnetico 2 polos 10A & 7 und & 7 unidad & 38.99 & 272.93 & E \\
interruptor termomagnetico 2 polos 32A & 1 und & 1 unidad & 37.90 & 37.90 & V \\
interruptor termomagnetico 2 polos 40A & 1 und & 1 unidad & 34.99 & 34.99 & E \\
interruptor termomagnetico 2 polos 63A & 1 und & 1 unidad & 35.45 & 35.45 & E \\
luminaria led plafon 18 w & 15 und & 15 unidad & 20.00 & 300.00 & E \\
conector abrazadera para varilla de puesta a tierra & 1 und & 1 unidad & 150.00 & 150.00 & E \\
varilla puesta a tierra cobre 5/8 2.4 m & 1 und & 1 unidad & 25.00 & 25.00 & E \\
tablero electrico 8 polos empotrable & 2 und & 2 unidad & 40.90 & 81.80 & V \\
tomacorriente doble 2p+t 16 a & 19 und & 19 unidad & 37.21 & 706.99 & E \\
tomacorriente especial 2p+t 20 a & 2 und & 2 unidad & 22.00 & 44.00 & E \\
tomacorriente especial 2p+t 32 a & 3 und & 3 unidad & 35.00 & 105.00 & E \\
\bottomrule\end{longtable}\normalsize

\subsection{Grupo 2: Sodimac/Maestro}
\scriptsize
\begin{longtable}{p{6.0cm}r r r r c}
\toprule Material & Metrado & Compra & P. comercial & Subtotal & Tipo \\ \midrule
\endfirsthead\toprule Material & Metrado & Compra & P. comercial & Subtotal & Tipo \\ \midrule\endhead
cinta aislante electrica & 3 und & 3 unidad & 6.72 & 20.16 & V \\
curva pvc sap electrica 20 mm & 16 und & 16 unidad & 2.60 & 41.60 & V \\
union pvc sap electrica 20 mm & 46 und & 46 unidad & 1.10 & 50.60 & E \\
union pvc sap electrica 40 mm & 4 und & 4 unidad & 2.50 & 10.00 & E \\
cable electrico thw 10 mm2 cobre & 119 m & 119 metro & 4.00 & 476.00 & E \\
cable electrico thw 10 mm2 cobre & 46 m & 46 metro & 3.50 & 161.00 & E \\
cable electrico thw 16 mm2 cobre & 40 m & 40 m & 19.90 & 796.00 & E \\
cable electrico thw 2.5 mm2 cobre & 262 m & 3 rollo & 209.90 & 629.70 & V \\
caja electrica estanca ip55 & 1 und & 1 unidad & 18.00 & 18.00 & E \\
caja electrica octogonal & 15 und & 15 paquete & 1.90 & 28.50 & V \\
caja electrica rectangular & 35 und & 35 paquete & 1.90 & 66.50 & V \\
caja registro puesta a tierra & 1 und & 1 unidad & 35.00 & 35.00 & E \\
tubo pvc sap electrico 20 mm & 152 m & 51 tubo & 6.90 & 351.90 & V \\
tubo pvc sap electrico 32 mm & 25 m & 9 tubo & 17.51 & 157.59 & V \\
tubo pvc sap electrico 40 mm & 13 m & 5 tubo & 12.00 & 60.00 & E \\
interruptor diferencial 2 polos 25A 30mA & 5 und & 5 unidad & 49.00 & 245.00 & V \\
interruptor diferencial 2 polos 40A 30mA & 2 und & 2 unidad & 65.90 & 131.80 & V \\
interruptor conmutado 10 a & 2 und & 2 unidad & 21.90 & 43.80 & E \\
interruptor simple 10 a & 9 und & 9 unidad & 43.90 & 395.10 & E \\
interruptor termomagnetico 2 polos 10A & 7 und & 7 unidad & 38.99 & 272.93 & V \\
interruptor termomagnetico 2 polos 32A & 1 und & 1 unidad & 31.90 & 31.90 & V \\
interruptor termomagnetico 2 polos 40A & 1 und & 1 unidad & 34.99 & 34.99 & V \\
interruptor termomagnetico 2 polos 63A & 1 und & 1 unidad & 29.90 & 29.90 & V \\
luminaria led plafon 18 w & 15 und & 15 unidad & 29.90 & 448.50 & E \\
conector abrazadera para varilla de puesta a tierra & 1 und & 1 unidad & 18.00 & 18.00 & V \\
varilla puesta a tierra cobre 5/8 2.4 m & 1 und & 1 unidad & 322.00 & 322.00 & E \\
tablero electrico 8 polos empotrable & 2 und & 2 unidad & 33.90 & 67.80 & V \\
tomacorriente doble 2p+t 16 a & 19 und & 19 unidad & 19.49 & 370.31 & V \\
tomacorriente especial 2p+t 20 a & 2 und & 2 unidad & 16.90 & 33.80 & E \\
tomacorriente especial 2p+t 32 a & 3 und & 3 unidad & 16.90 & 50.70 & E \\
\bottomrule\end{longtable}\normalsize

\subsection{Grupo 3: Mercado Libre Peru}
\scriptsize
\begin{longtable}{p{6.0cm}r r r r c}
\toprule Material & Metrado & Compra & P. comercial & Subtotal & Tipo \\ \midrule
\endfirsthead\toprule Material & Metrado & Compra & P. comercial & Subtotal & Tipo \\ \midrule\endhead
cinta aislante electrica & 3 und & 3 unidad & 5.71 & 17.13 & E \\
curva pvc sap electrica 20 mm & 16 und & 16 unidad & 2.48 & 39.60 & E \\
union pvc sap electrica 20 mm & 46 und & 46 unidad & 1.10 & 50.60 & E \\
union pvc sap electrica 40 mm & 4 und & 4 unidad & 2.50 & 10.00 & E \\
cable electrico thw 10 mm2 cobre & 119 m & 119 metro & 4.00 & 476.00 & E \\
cable electrico thw 10 mm2 cobre & 46 m & 46 metro & 3.50 & 161.00 & E \\
cable electrico thw 16 mm2 cobre & 40 m & 40 metro & 6.00 & 240.00 & E \\
cable electrico thw 2.5 mm2 cobre & 262 m & 3 rollo & 209.90 & 629.70 & E \\
caja electrica estanca ip55 & 1 und & 1 unidad & 18.00 & 18.00 & E \\
caja electrica octogonal & 15 und & 15 paquete & 1.90 & 28.50 & E \\
caja electrica rectangular & 35 und & 35 paquete & 1.94 & 67.90 & E \\
caja registro puesta a tierra & 1 und & 1 unidad & 35.00 & 35.00 & E \\
tubo pvc sap electrico 20 mm & 152 m & 51 tubo & 8.40 & 428.40 & E \\
tubo pvc sap electrico 32 mm & 25 m & 9 tubo & 18.31 & 164.76 & E \\
tubo pvc sap electrico 40 mm & 13 m & 5 tubo & 12.00 & 60.00 & E \\
interruptor diferencial 2 polos 25A 30mA & 5 und & 5 unidad & 49.49 & 247.47 & E \\
interruptor diferencial 2 polos 40A 30mA & 2 und & 2 unidad & 65.90 & 131.80 & E \\
interruptor conmutado 10 a & 2 und & 2 unidad & 18.00 & 36.00 & E \\
interruptor simple 10 a & 9 und & 9 unidad & 12.00 & 108.00 & E \\
interruptor termomagnetico 2 polos 10A & 7 und & 7 unidad & 38.99 & 272.93 & E \\
interruptor termomagnetico 2 polos 32A & 1 und & 1 unidad & 35.17 & 35.17 & E \\
interruptor termomagnetico 2 polos 40A & 1 und & 1 unidad & 34.99 & 34.99 & E \\
interruptor termomagnetico 2 polos 63A & 1 und & 1 unidad & 35.45 & 35.45 & E \\
luminaria led plafon 18 w & 15 und & 15 unidad & 25.00 & 375.00 & E \\
conector abrazadera para varilla de puesta a tierra & 1 und & 1 unidad & 20.45 & 20.45 & E \\
varilla puesta a tierra cobre 5/8 2.4 m & 1 und & 1 unidad & 25.00 & 25.00 & E \\
tablero electrico 8 polos empotrable & 2 und & 2 unidad & 50.01 & 100.02 & E \\
tomacorriente doble 2p+t 16 a & 19 und & 19 unidad & 37.21 & 706.99 & E \\
tomacorriente especial 2p+t 20 a & 2 und & 2 unidad & 22.00 & 44.00 & E \\
tomacorriente especial 2p+t 32 a & 3 und & 3 unidad & 35.00 & 105.00 & E \\
\bottomrule\end{longtable}\normalsize

\section{Comparativa economica}
\scriptsize
\begin{longtable}{p{4.5cm}rrrrp{2.5cm}r}
\toprule Material & Promart & Sod./Mae. & M. Libre & Mejor & Proveedor & Recomendado \\ \midrule
\endfirsthead\toprule Material & Promart & Sod./Mae. & M. Libre & Mejor & Proveedor & Recomendado \\ \midrule\endhead
cable electrico thw 10 mm2 cobre & 476.00 & 476.00 & 476.00 & Promart & 476.00 \\
cable electrico thw 10 mm2 cobre & 161.00 & 161.00 & 161.00 & Promart & 161.00 \\
cable electrico thw 16 mm2 cobre & 240.00 & 796.00 & 240.00 & Promart & 240.00 \\
cable electrico thw 2.5 mm2 cobre & 629.70 & 629.70 & 629.70 & Sodimac/Maestro & 629.70 \\
caja electrica estanca ip55 & 18.00 & 18.00 & 18.00 & Promart & 18.00 \\
caja electrica octogonal & 180.00 & 28.50 & 28.50 & Sodimac/Maestro & 28.50 \\
caja electrica rectangular & 67.90 & 66.50 & 67.90 & Sodimac/Maestro & 66.50 \\
caja registro puesta a tierra & 35.00 & 35.00 & 35.00 & Promart & 35.00 \\
cinta aislante electrica & 6.00 & 20.16 & 17.13 & Promart & 6.00 \\
conector abrazadera para varilla de puesta a tierra & 150.00 & 18.00 & 20.45 & Sodimac/Maestro & 18.00 \\
curva pvc sap electrica 20 mm & 35.20 & 41.60 & 39.60 & Promart & 35.20 \\
interruptor conmutado 10 a & 27.80 & 43.80 & 36.00 & Promart & 27.80 \\
interruptor diferencial 2 polos 25A 30mA & 247.47 & 245.00 & 247.47 & Sodimac/Maestro & 245.00 \\
interruptor diferencial 2 polos 40A 30mA & 131.80 & 131.80 & 131.80 & Sodimac/Maestro & 131.80 \\
interruptor simple 10 a & 53.10 & 395.10 & 108.00 & Promart & 53.10 \\
interruptor termomagnetico 2 polos 10A & 272.93 & 272.93 & 272.93 & Sodimac/Maestro & 272.93 \\
interruptor termomagnetico 2 polos 32A & 37.90 & 31.90 & 35.17 & Sodimac/Maestro & 31.90 \\
interruptor termomagnetico 2 polos 40A & 34.99 & 34.99 & 34.99 & Sodimac/Maestro & 34.99 \\
interruptor termomagnetico 2 polos 63A & 35.45 & 29.90 & 35.45 & Sodimac/Maestro & 29.90 \\
luminaria led plafon 18 w & 300.00 & 448.50 & 375.00 & Promart & 300.00 \\
tablero electrico 8 polos empotrable & 81.80 & 67.80 & 100.02 & Sodimac/Maestro & 67.80 \\
tomacorriente doble 2p+t 16 a & 706.99 & 370.31 & 706.99 & Sodimac/Maestro & 370.31 \\
tomacorriente especial 2p+t 20 a & 44.00 & 33.80 & 44.00 & Sodimac/Maestro & 33.80 \\
tomacorriente especial 2p+t 32 a & 105.00 & 50.70 & 105.00 & Sodimac/Maestro & 50.70 \\
tubo pvc sap electrico 20 mm & 300.90 & 351.90 & 428.40 & Promart & 300.90 \\
tubo pvc sap electrico 32 mm & 179.10 & 157.59 & 164.76 & Sodimac/Maestro & 157.59 \\
tubo pvc sap electrico 40 mm & 60.00 & 60.00 & 60.00 & Promart & 60.00 \\
union pvc sap electrica 20 mm & 50.60 & 50.60 & 50.60 & Promart & 50.60 \\
union pvc sap electrica 40 mm & 10.00 & 10.00 & 10.00 & Promart & 10.00 \\
varilla puesta a tierra cobre 5/8 2.4 m & 25.00 & 322.00 & 25.00 & Promart & 25.00 \\
\bottomrule\end{longtable}\normalsize

\section{Presupuesto recomendado}
\scriptsize
\begin{longtable}{p{5.5cm}p{2.5cm}r r p{3.0cm}}
\toprule Material & Proveedor & Compra & Subtotal & Tipo de precio \\ \midrule
\endfirsthead\toprule Material & Proveedor & Compra & Subtotal & Tipo de precio \\ \midrule\endhead
cinta aislante electrica & Promart & 3 unidad & 6.00 & Verificado (Tienda) \\
curva pvc sap electrica 20 mm & Promart & 16 unidad & 35.20 & Verificado (Tienda) \\
union pvc sap electrica 20 mm & Promart & 46 unidad & 50.60 & Verificado (Tienda) \\
union pvc sap electrica 40 mm & Promart & 4 unidad & 10.00 & Estimado (Referencia) \\
cable electrico thw 10 mm2 cobre & Promart & 119 metro & 476.00 & Estimado (Referencia) \\
cable electrico thw 10 mm2 cobre & Promart & 46 metro & 161.00 & Estimado (Referencia) \\
cable electrico thw 16 mm2 cobre & Promart & 40 metro & 240.00 & Estimado (Referencia) \\
cable electrico thw 2.5 mm2 cobre & Sodimac/Maestro & 3 rollo & 629.70 & Verificado (Tienda) \\
caja electrica estanca ip55 & Promart & 1 unidad & 18.00 & Estimado (Referencia) \\
caja electrica octogonal & Sodimac/Maestro & 15 paquete & 28.50 & Verificado (Tienda) \\
caja electrica rectangular & Sodimac/Maestro & 35 paquete & 66.50 & Verificado (Tienda) \\
caja registro puesta a tierra & Promart & 1 unidad & 35.00 & Estimado (Referencia) \\
tubo pvc sap electrico 20 mm & Promart & 51 tubo & 300.90 & Verificado (Tienda) \\
tubo pvc sap electrico 32 mm & Sodimac/Maestro & 9 tubo & 157.59 & Verificado (Tienda) \\
tubo pvc sap electrico 40 mm & Promart & 5 tubo & 60.00 & Estimado (Referencia) \\
interruptor diferencial 2 polos 25A 30mA & Sodimac/Maestro & 5 unidad & 245.00 & Verificado (Tienda) \\
interruptor diferencial 2 polos 40A 30mA & Sodimac/Maestro & 2 unidad & 131.80 & Verificado (Tienda) \\
interruptor conmutado 10 a & Promart & 2 unidad & 27.80 & Estimado (Equivalente) \\
interruptor simple 10 a & Promart & 9 unidad & 53.10 & Estimado (Equivalente) \\
interruptor termomagnetico 2 polos 10A & Sodimac/Maestro & 7 unidad & 272.93 & Verificado (Tienda) \\
interruptor termomagnetico 2 polos 32A & Sodimac/Maestro & 1 unidad & 31.90 & Verificado (Tienda) \\
interruptor termomagnetico 2 polos 40A & Sodimac/Maestro & 1 unidad & 34.99 & Verificado (Tienda) \\
interruptor termomagnetico 2 polos 63A & Sodimac/Maestro & 1 unidad & 29.90 & Verificado (Tienda) \\
luminaria led plafon 18 w & Promart & 15 unidad & 300.00 & Estimado (Equivalente) \\
conector abrazadera para varilla de puesta a tierra & Sodimac/Maestro & 1 unidad & 18.00 & Verificado (Tienda) \\
varilla puesta a tierra cobre 5/8 2.4 m & Promart & 1 unidad & 25.00 & Estimado (Referencia) \\
tablero electrico 8 polos empotrable & Sodimac/Maestro & 2 unidad & 67.80 & Verificado (Tienda) \\
tomacorriente doble 2p+t 16 a & Sodimac/Maestro & 19 unidad & 370.31 & Verificado (Tienda) \\
tomacorriente especial 2p+t 20 a & Sodimac/Maestro & 2 unidad & 33.80 & Estimado (Equivalente) \\
tomacorriente especial 2p+t 32 a & Sodimac/Maestro & 3 unidad & 50.70 & Estimado (Equivalente) \\

\midrule\multicolumn{3}{r}{\textbf{TOTAL}} & \textbf{S/ 3,968.02} & \\ \bottomrule\end{longtable}\normalsize

% Requiere: booktabs, longtable, array, xcolor, hyperref.
\section{Metrados de materiales electricos}
Se validaron 68 partidas usando los DXF vigentes, sus JSON documentales y los calculos previos del BOM.
\scriptsize
\begin{longtable}{p{1.7cm}p{7.5cm}r c p{2.2cm}}
\toprule Codigo & Descripcion & Cant. & Und. & Circuito \\ \midrule
\endfirsthead\toprule Codigo & Descripcion & Cant. & Und. & Circuito \\ \midrule\endhead
BOM-001 & Cable TW THW 16 mm2 (alimentador trifasico) & 40 & m & ALIM-TG \\
BOM-002 & Tubo PVC SAP 40 mm (alimentador) & 13 & m & ALIM-TG \\
BOM-003 & ITM 2P-63A & 1 & und & ALIM-TG \\
BOM-004 & Tablero general para 8 circuitos & 1 & und & GENERAL \\
BOM-005 & Interruptor diferencial 2P-40A-30mA & 1 & und & GENERAL \\
BOM-006 & Cable TW THW 2.5 mm2 - C1 & 33 & m & C1 \\
BOM-007 & Tubo PVC SAP 20 mm - C1 & 16 & m & C1 \\
BOM-008 & ITM 2P-10A - C1 & 1 & und & C1 \\
BOM-010 & Cable TW THW 2.5 mm2 - C2 & 40 & m & C2 \\
BOM-011 & Tubo PVC SAP 20 mm - C2 & 20 & m & C2 \\
BOM-012 & ITM 2P-10A - C2 & 1 & und & C2 \\
BOM-013 & Diferencial 2P-25A-30mA - C2 & 1 & und & C2 \\
BOM-015 & Cable TW THW 2.5 mm2 - C3 & 22 & m & C3 \\
BOM-016 & Tubo PVC SAP 20 mm - C3 & 11 & m & C3 \\
BOM-017 & ITM 2P-10A - C3 & 1 & und & C3 \\
BOM-018 & Diferencial 2P-25A-30mA - C3 & 1 & und & C3 \\
BOM-020 & Cable TW THW 2.5 mm2 - C4 & 48 & m & C4 \\
BOM-021 & Tubo PVC SAP 20 mm - C4 & 24 & m & C4 \\
BOM-022 & ITM 2P-10A - C4 & 1 & und & C4 \\
BOM-024 & Cable TW THW 2.5 mm2 - C5 & 55 & m & C5 \\
BOM-025 & Tubo PVC SAP 20 mm - C5 & 28 & m & C5 \\
BOM-026 & ITM 2P-10A - C5 & 1 & und & C5 \\
BOM-027 & Diferencial 2P-25A-30mA - C5 & 1 & und & C5 \\
BOM-029 & Cable TW THW 10 mm2 - C6 & 44 & m & C6 \\
BOM-030 & Tubo PVC SAP 20 mm - C6 & 22 & m & C6 \\
BOM-031 & ITM 2P-32A - C6 & 1 & und & C6 \\
BOM-032 & Diferencial 2P-40A-30mA - C6 & 1 & und & C6 \\
BOM-034 & Cable TW THW 2.5 mm2 - C7 & 33 & m & C7 \\
BOM-035 & Tubo PVC SAP 20 mm - C7 & 16 & m & C7 \\
BOM-036 & ITM 2P-10A - C7 & 1 & und & C7 \\
BOM-037 & Diferencial 2P-25A-30mA - C7 & 1 & und & C7 \\
BOM-039 & Cable TW THW 2.5 mm2 - C8 & 31 & m & C8 \\
BOM-040 & Tubo PVC SAP 20 mm - C8 & 15 & m & C8 \\
BOM-041 & ITM 2P-10A - C8 & 1 & und & C8 \\
BOM-042 & Diferencial 2P-25A-30mA - C8 & 1 & und & C8 \\
BOM-044 & Cable de tierra TW 10 mm2 (verde/amarillo) & 46 & m & SPAT \\
BOM-045 & Varilla de puesta a tierra 5/8 x 2.4 m & 1 & und & SPAT \\
BOM-046 & Cinta de aislar electrica & 3 & und & GENERAL \\
BOM-047 & Caja octogonal para luminaria - C1 & 4 & und & C1 \\
BOM-048 & Caja rectangular para interruptor - C1 & 4 & und & C1 \\
BOM-049 & Caja rectangular para tomacorriente - C2 & 7 & und & C2 \\
BOM-050 & Caja rectangular para tomacorriente especial - C3 & 1 & und & C3 \\
BOM-051 & Caja octogonal para luminaria - C4 & 11 & und & C4 \\
BOM-052 & Caja rectangular para interruptor - C4 & 7 & und & C4 \\
BOM-053 & Caja rectangular para tomacorriente - C5 & 12 & und & C5 \\
BOM-054 & Caja rectangular para salida especial - C6 & 3 & und & C6 \\
BOM-055 & Caja rectangular para tomacorriente especial - C7 & 1 & und & C7 \\
BOM-056 & Caja estanca IP55 para conexion de bomba - C8 & 1 & und & C8 \\
BOM-057 & Luminaria LED plafon 18 W - C1 & 4 & und & C1 \\
BOM-058 & Luminaria LED plafon 18 W - C4 & 11 & und & C4 \\
BOM-059 & Interruptor simple 10 A - C1 & 3 & und & C1 \\
BOM-060 & Interruptor conmutado 10 A - C1 & 1 & und & C1 \\
BOM-061 & Interruptor simple 10 A - C4 & 6 & und & C4 \\
BOM-062 & Interruptor conmutado 10 A - C4 & 1 & und & C4 \\
BOM-063 & Tomacorriente doble 2P+T 16 A - C2 & 7 & und & C2 \\
BOM-064 & Tomacorriente doble 2P+T 16 A - C5 & 12 & und & C5 \\
BOM-065 & Tomacorriente especial 2P+T 20 A - C3 & 1 & und & C3 \\
BOM-066 & Tomacorriente especial 2P+T 32 A - C6 & 3 & und & C6 \\
BOM-067 & Tomacorriente especial 2P+T 20 A - C7 & 1 & und & C7 \\
BOM-068 & Tablero T2 para 8 circuitos & 1 & und & T2 \\
BOM-069 & Caja de registro para puesta a tierra & 1 & und & SPAT \\
BOM-070 & Conector abrazadera para varilla de puesta a tierra & 1 & und & SPAT \\
BOM-071 & Union PVC SAP 20 mm & 46 & und & C1-C8 \\
BOM-072 & Curva PVC SAP 20 mm & 16 & und & C1-C8 \\
BOM-073 & Union PVC SAP 40 mm & 4 & und & ALIM-TG \\
BOM-074 & Cable TW THW 10 mm2 - alimentador TG a T2 & 75 & m & ALIM-T2 \\
BOM-075 & Tubo PVC SAP 32 mm - alimentador TG a T2 & 25 & m & ALIM-T2 \\
BOM-076 & ITM 2P-40A - alimentador T2 & 1 & und & ALIM-T2 \\
\bottomrule\end{longtable}\normalsize

\section{Cotizacion por proveedor}
\subsection{Grupo 1: Promart}
\scriptsize
\begin{longtable}{p{6.0cm}r r r r c}
\toprule Material & Metrado & Compra & P. comercial & Subtotal & Tipo \\ \midrule
\endfirsthead\toprule Material & Metrado & Compra & P. comercial & Subtotal & Tipo \\ \midrule\endhead
cinta aislante electrica & 3 und & 3 unidad & 2.00 & 6.00 & V \\
curva pvc sap electrica 20 mm & 16 und & 16 unidad & 2.20 & 35.20 & V \\
union pvc sap electrica 20 mm & 46 und & 46 unidad & 1.10 & 50.60 & V \\
union pvc sap electrica 40 mm & 4 und & 4 unidad & 2.50 & 10.00 & E \\
cable electrico thw 10 mm2 cobre & 119 m & 119 metro & 4.00 & 476.00 & E \\
cable electrico thw 10 mm2 cobre & 46 m & 46 metro & 3.50 & 161.00 & E \\
cable electrico thw 16 mm2 cobre & 40 m & 40 metro & 6.00 & 240.00 & E \\
cable electrico thw 2.5 mm2 cobre & 262 m & 3 rollo & 209.90 & 629.70 & E \\
caja electrica estanca ip55 & 1 und & 1 unidad & 18.00 & 18.00 & E \\
caja electrica octogonal & 15 und & 15 unidad & 12.00 & 180.00 & E \\
caja electrica rectangular & 35 und & 35 paquete & 1.94 & 67.90 & E \\
caja registro puesta a tierra & 1 und & 1 unidad & 35.00 & 35.00 & E \\
tubo pvc sap electrico 20 mm & 152 m & 51 tubo & 5.90 & 300.90 & V \\
tubo pvc sap electrico 32 mm & 25 m & 9 tubo & 19.90 & 179.10 & V \\
tubo pvc sap electrico 40 mm & 13 m & 5 tubo & 12.00 & 60.00 & E \\
interruptor diferencial 2 polos 25A 30mA & 5 und & 5 unidad & 49.49 & 247.47 & E \\
interruptor diferencial 2 polos 40A 30mA & 2 und & 2 unidad & 65.90 & 131.80 & E \\
interruptor conmutado 10 a & 2 und & 2 unidad & 13.90 & 27.80 & E \\
interruptor simple 10 a & 9 und & 9 unidad & 5.90 & 53.10 & E \\
interruptor termomagnetico 2 polos 10A & 7 und & 7 unidad & 38.99 & 272.93 & E \\
interruptor termomagnetico 2 polos 32A & 1 und & 1 unidad & 37.90 & 37.90 & V \\
interruptor termomagnetico 2 polos 40A & 1 und & 1 unidad & 34.99 & 34.99 & E \\
interruptor termomagnetico 2 polos 63A & 1 und & 1 unidad & 35.45 & 35.45 & E \\
luminaria led plafon 18 w & 15 und & 15 unidad & 20.00 & 300.00 & E \\
conector abrazadera para varilla de puesta a tierra & 1 und & 1 unidad & 150.00 & 150.00 & E \\
varilla puesta a tierra cobre 5/8 2.4 m & 1 und & 1 unidad & 25.00 & 25.00 & E \\
tablero electrico 8 polos empotrable & 2 und & 2 unidad & 40.90 & 81.80 & V \\
tomacorriente doble 2p+t 16 a & 19 und & 19 unidad & 37.21 & 706.99 & E \\
tomacorriente especial 2p+t 20 a & 2 und & 2 unidad & 22.00 & 44.00 & E \\
tomacorriente especial 2p+t 32 a & 3 und & 3 unidad & 35.00 & 105.00 & E \\
\bottomrule\end{longtable}\normalsize

\subsection{Grupo 2: Sodimac/Maestro}
\scriptsize
\begin{longtable}{p{6.0cm}r r r r c}
\toprule Material & Metrado & Compra & P. comercial & Subtotal & Tipo \\ \midrule
\endfirsthead\toprule Material & Metrado & Compra & P. comercial & Subtotal & Tipo \\ \midrule\endhead
cinta aislante electrica & 3 und & 3 unidad & 6.72 & 20.16 & V \\
curva pvc sap electrica 20 mm & 16 und & 16 unidad & 2.60 & 41.60 & V \\
union pvc sap electrica 20 mm & 46 und & 46 unidad & 1.10 & 50.60 & E \\
union pvc sap electrica 40 mm & 4 und & 4 unidad & 2.50 & 10.00 & E \\
cable electrico thw 10 mm2 cobre & 119 m & 119 metro & 4.00 & 476.00 & E \\
cable electrico thw 10 mm2 cobre & 46 m & 46 metro & 3.50 & 161.00 & E \\
cable electrico thw 16 mm2 cobre & 40 m & 40 m & 19.90 & 796.00 & E \\
cable electrico thw 2.5 mm2 cobre & 262 m & 3 rollo & 209.90 & 629.70 & V \\
caja electrica estanca ip55 & 1 und & 1 unidad & 18.00 & 18.00 & E \\
caja electrica octogonal & 15 und & 15 paquete & 1.90 & 28.50 & V \\
caja electrica rectangular & 35 und & 35 paquete & 1.90 & 66.50 & V \\
caja registro puesta a tierra & 1 und & 1 unidad & 35.00 & 35.00 & E \\
tubo pvc sap electrico 20 mm & 152 m & 51 tubo & 6.90 & 351.90 & V \\
tubo pvc sap electrico 32 mm & 25 m & 9 tubo & 17.51 & 157.59 & V \\
tubo pvc sap electrico 40 mm & 13 m & 5 tubo & 12.00 & 60.00 & E \\
interruptor diferencial 2 polos 25A 30mA & 5 und & 5 unidad & 49.00 & 245.00 & V \\
interruptor diferencial 2 polos 40A 30mA & 2 und & 2 unidad & 65.90 & 131.80 & V \\
interruptor conmutado 10 a & 2 und & 2 unidad & 21.90 & 43.80 & E \\
interruptor simple 10 a & 9 und & 9 unidad & 43.90 & 395.10 & E \\
interruptor termomagnetico 2 polos 10A & 7 und & 7 unidad & 38.99 & 272.93 & V \\
interruptor termomagnetico 2 polos 32A & 1 und & 1 unidad & 31.90 & 31.90 & V \\
interruptor termomagnetico 2 polos 40A & 1 und & 1 unidad & 34.99 & 34.99 & V \\
interruptor termomagnetico 2 polos 63A & 1 und & 1 unidad & 29.90 & 29.90 & V \\
luminaria led plafon 18 w & 15 und & 15 unidad & 29.90 & 448.50 & E \\
conector abrazadera para varilla de puesta a tierra & 1 und & 1 unidad & 18.00 & 18.00 & V \\
varilla puesta a tierra cobre 5/8 2.4 m & 1 und & 1 unidad & 322.00 & 322.00 & E \\
tablero electrico 8 polos empotrable & 2 und & 2 unidad & 33.90 & 67.80 & V \\
tomacorriente doble 2p+t 16 a & 19 und & 19 unidad & 19.49 & 370.31 & V \\
tomacorriente especial 2p+t 20 a & 2 und & 2 unidad & 16.90 & 33.80 & E \\
tomacorriente especial 2p+t 32 a & 3 und & 3 unidad & 16.90 & 50.70 & E \\
\bottomrule\end{longtable}\normalsize

\subsection{Grupo 3: Mercado Libre Peru}
\scriptsize
\begin{longtable}{p{6.0cm}r r r r c}
\toprule Material & Metrado & Compra & P. comercial & Subtotal & Tipo \\ \midrule
\endfirsthead\toprule Material & Metrado & Compra & P. comercial & Subtotal & Tipo \\ \midrule\endhead
cinta aislante electrica & 3 und & 3 unidad & 5.71 & 17.13 & E \\
curva pvc sap electrica 20 mm & 16 und & 16 unidad & 2.48 & 39.60 & E \\
union pvc sap electrica 20 mm & 46 und & 46 unidad & 1.10 & 50.60 & E \\
union pvc sap electrica 40 mm & 4 und & 4 unidad & 2.50 & 10.00 & E \\
cable electrico thw 10 mm2 cobre & 119 m & 119 metro & 4.00 & 476.00 & E \\
cable electrico thw 10 mm2 cobre & 46 m & 46 metro & 3.50 & 161.00 & E \\
cable electrico thw 16 mm2 cobre & 40 m & 40 metro & 6.00 & 240.00 & E \\
cable electrico thw 2.5 mm2 cobre & 262 m & 3 rollo & 209.90 & 629.70 & E \\
caja electrica estanca ip55 & 1 und & 1 unidad & 18.00 & 18.00 & E \\
caja electrica octogonal & 15 und & 15 paquete & 1.90 & 28.50 & E \\
caja electrica rectangular & 35 und & 35 paquete & 1.94 & 67.90 & E \\
caja registro puesta a tierra & 1 und & 1 unidad & 35.00 & 35.00 & E \\
tubo pvc sap electrico 20 mm & 152 m & 51 tubo & 8.40 & 428.40 & E \\
tubo pvc sap electrico 32 mm & 25 m & 9 tubo & 18.31 & 164.76 & E \\
tubo pvc sap electrico 40 mm & 13 m & 5 tubo & 12.00 & 60.00 & E \\
interruptor diferencial 2 polos 25A 30mA & 5 und & 5 unidad & 49.49 & 247.47 & E \\
interruptor diferencial 2 polos 40A 30mA & 2 und & 2 unidad & 65.90 & 131.80 & E \\
interruptor conmutado 10 a & 2 und & 2 unidad & 18.00 & 36.00 & E \\
interruptor simple 10 a & 9 und & 9 unidad & 12.00 & 108.00 & E \\
interruptor termomagnetico 2 polos 10A & 7 und & 7 unidad & 38.99 & 272.93 & E \\
interruptor termomagnetico 2 polos 32A & 1 und & 1 unidad & 35.17 & 35.17 & E \\
interruptor termomagnetico 2 polos 40A & 1 und & 1 unidad & 34.99 & 34.99 & E \\
interruptor termomagnetico 2 polos 63A & 1 und & 1 unidad & 35.45 & 35.45 & E \\
luminaria led plafon 18 w & 15 und & 15 unidad & 25.00 & 375.00 & E \\
conector abrazadera para varilla de puesta a tierra & 1 und & 1 unidad & 20.45 & 20.45 & E \\
varilla puesta a tierra cobre 5/8 2.4 m & 1 und & 1 unidad & 25.00 & 25.00 & E \\
tablero electrico 8 polos empotrable & 2 und & 2 unidad & 50.01 & 100.02 & E \\
tomacorriente doble 2p+t 16 a & 19 und & 19 unidad & 37.21 & 706.99 & E \\
tomacorriente especial 2p+t 20 a & 2 und & 2 unidad & 22.00 & 44.00 & E \\
tomacorriente especial 2p+t 32 a & 3 und & 3 unidad & 35.00 & 105.00 & E \\
\bottomrule\end{longtable}\normalsize

\section{Comparativa economica}
\scriptsize
\begin{longtable}{p{4.5cm}rrrrp{2.5cm}r}
\toprule Material & Promart & Sod./Mae. & M. Libre & Mejor & Proveedor & Recomendado \\ \midrule
\endfirsthead\toprule Material & Promart & Sod./Mae. & M. Libre & Mejor & Proveedor & Recomendado \\ \midrule\endhead
cable electrico thw 10 mm2 cobre & 476.00 & 476.00 & 476.00 & Promart & 476.00 \\
cable electrico thw 10 mm2 cobre & 161.00 & 161.00 & 161.00 & Promart & 161.00 \\
cable electrico thw 16 mm2 cobre & 240.00 & 796.00 & 240.00 & Promart & 240.00 \\
cable electrico thw 2.5 mm2 cobre & 629.70 & 629.70 & 629.70 & Sodimac/Maestro & 629.70 \\
caja electrica estanca ip55 & 18.00 & 18.00 & 18.00 & Promart & 18.00 \\
caja electrica octogonal & 180.00 & 28.50 & 28.50 & Sodimac/Maestro & 28.50 \\
caja electrica rectangular & 67.90 & 66.50 & 67.90 & Sodimac/Maestro & 66.50 \\
caja registro puesta a tierra & 35.00 & 35.00 & 35.00 & Promart & 35.00 \\
cinta aislante electrica & 6.00 & 20.16 & 17.13 & Promart & 6.00 \\
conector abrazadera para varilla de puesta a tierra & 150.00 & 18.00 & 20.45 & Sodimac/Maestro & 18.00 \\
curva pvc sap electrica 20 mm & 35.20 & 41.60 & 39.60 & Promart & 35.20 \\
interruptor conmutado 10 a & 27.80 & 43.80 & 36.00 & Promart & 27.80 \\
interruptor diferencial 2 polos 25A 30mA & 247.47 & 245.00 & 247.47 & Sodimac/Maestro & 245.00 \\
interruptor diferencial 2 polos 40A 30mA & 131.80 & 131.80 & 131.80 & Sodimac/Maestro & 131.80 \\
interruptor simple 10 a & 53.10 & 395.10 & 108.00 & Promart & 53.10 \\
interruptor termomagnetico 2 polos 10A & 272.93 & 272.93 & 272.93 & Sodimac/Maestro & 272.93 \\
interruptor termomagnetico 2 polos 32A & 37.90 & 31.90 & 35.17 & Sodimac/Maestro & 31.90 \\
interruptor termomagnetico 2 polos 40A & 34.99 & 34.99 & 34.99 & Sodimac/Maestro & 34.99 \\
interruptor termomagnetico 2 polos 63A & 35.45 & 29.90 & 35.45 & Sodimac/Maestro & 29.90 \\
luminaria led plafon 18 w & 300.00 & 448.50 & 375.00 & Promart & 300.00 \\
tablero electrico 8 polos empotrable & 81.80 & 67.80 & 100.02 & Sodimac/Maestro & 67.80 \\
tomacorriente doble 2p+t 16 a & 706.99 & 370.31 & 706.99 & Sodimac/Maestro & 370.31 \\
tomacorriente especial 2p+t 20 a & 44.00 & 33.80 & 44.00 & Sodimac/Maestro & 33.80 \\
tomacorriente especial 2p+t 32 a & 105.00 & 50.70 & 105.00 & Sodimac/Maestro & 50.70 \\
tubo pvc sap electrico 20 mm & 300.90 & 351.90 & 428.40 & Promart & 300.90 \\
tubo pvc sap electrico 32 mm & 179.10 & 157.59 & 164.76 & Sodimac/Maestro & 157.59 \\
tubo pvc sap electrico 40 mm & 60.00 & 60.00 & 60.00 & Promart & 60.00 \\
union pvc sap electrica 20 mm & 50.60 & 50.60 & 50.60 & Promart & 50.60 \\
union pvc sap electrica 40 mm & 10.00 & 10.00 & 10.00 & Promart & 10.00 \\
varilla puesta a tierra cobre 5/8 2.4 m & 25.00 & 322.00 & 25.00 & Promart & 25.00 \\
\bottomrule\end{longtable}\normalsize

\section{Presupuesto recomendado}
\scriptsize
\begin{longtable}{p{5.5cm}p{2.5cm}r r p{3.0cm}}
\toprule Material & Proveedor & Compra & Subtotal & Tipo de precio \\ \midrule
\endfirsthead\toprule Material & Proveedor & Compra & Subtotal & Tipo de precio \\ \midrule\endhead
cinta aislante electrica & Promart & 3 unidad & 6.00 & Verificado (Tienda) \\
curva pvc sap electrica 20 mm & Promart & 16 unidad & 35.20 & Verificado (Tienda) \\
union pvc sap electrica 20 mm & Promart & 46 unidad & 50.60 & Verificado (Tienda) \\
union pvc sap electrica 40 mm & Promart & 4 unidad & 10.00 & Estimado (Referencia) \\
cable electrico thw 10 mm2 cobre & Promart & 119 metro & 476.00 & Estimado (Referencia) \\
cable electrico thw 10 mm2 cobre & Promart & 46 metro & 161.00 & Estimado (Referencia) \\
cable electrico thw 16 mm2 cobre & Promart & 40 metro & 240.00 & Estimado (Referencia) \\
cable electrico thw 2.5 mm2 cobre & Sodimac/Maestro & 3 rollo & 629.70 & Verificado (Tienda) \\
caja electrica estanca ip55 & Promart & 1 unidad & 18.00 & Estimado (Referencia) \\
caja electrica octogonal & Sodimac/Maestro & 15 paquete & 28.50 & Verificado (Tienda) \\
caja electrica rectangular & Sodimac/Maestro & 35 paquete & 66.50 & Verificado (Tienda) \\
caja registro puesta a tierra & Promart & 1 unidad & 35.00 & Estimado (Referencia) \\
tubo pvc sap electrico 20 mm & Promart & 51 tubo & 300.90 & Verificado (Tienda) \\
tubo pvc sap electrico 32 mm & Sodimac/Maestro & 9 tubo & 157.59 & Verificado (Tienda) \\
tubo pvc sap electrico 40 mm & Promart & 5 tubo & 60.00 & Estimado (Referencia) \\
interruptor diferencial 2 polos 25A 30mA & Sodimac/Maestro & 5 unidad & 245.00 & Verificado (Tienda) \\
interruptor diferencial 2 polos 40A 30mA & Sodimac/Maestro & 2 unidad & 131.80 & Verificado (Tienda) \\
interruptor conmutado 10 a & Promart & 2 unidad & 27.80 & Estimado (Equivalente) \\
interruptor simple 10 a & Promart & 9 unidad & 53.10 & Estimado (Equivalente) \\
interruptor termomagnetico 2 polos 10A & Sodimac/Maestro & 7 unidad & 272.93 & Verificado (Tienda) \\
interruptor termomagnetico 2 polos 32A & Sodimac/Maestro & 1 unidad & 31.90 & Verificado (Tienda) \\
interruptor termomagnetico 2 polos 40A & Sodimac/Maestro & 1 unidad & 34.99 & Verificado (Tienda) \\
interruptor termomagnetico 2 polos 63A & Sodimac/Maestro & 1 unidad & 29.90 & Verificado (Tienda) \\
luminaria led plafon 18 w & Promart & 15 unidad & 300.00 & Estimado (Equivalente) \\
conector abrazadera para varilla de puesta a tierra & Sodimac/Maestro & 1 unidad & 18.00 & Verificado (Tienda) \\
varilla puesta a tierra cobre 5/8 2.4 m & Promart & 1 unidad & 25.00 & Estimado (Referencia) \\
tablero electrico 8 polos empotrable & Sodimac/Maestro & 2 unidad & 67.80 & Verificado (Tienda) \\
tomacorriente doble 2p+t 16 a & Sodimac/Maestro & 19 unidad & 370.31 & Verificado (Tienda) \\
tomacorriente especial 2p+t 20 a & Sodimac/Maestro & 2 unidad & 33.80 & Estimado (Equivalente) \\
tomacorriente especial 2p+t 32 a & Sodimac/Maestro & 3 unidad & 50.70 & Estimado (Equivalente) \\

\midrule\multicolumn{3}{r}{\textbf{TOTAL}} & \textbf{S/ 3,968.02} & \\ \bottomrule\end{longtable}\normalsize


\end{document}
